\documentclass[10.5pt,a5paper]{article}

\usepackage{fontspec}
\setmainfont{DejaVu Serif}
\setsansfont{DejaVu Sans}
\setmonofont{DejaVu Sans Mono}
\usepackage[russian]{babel}
\usepackage{geometry}
\geometry{left=13mm,right=13mm,top=14mm,bottom=16mm}
\usepackage{microtype}
\usepackage{array,tabularx,booktabs}
\usepackage{enumitem}
\usepackage{xcolor}
\usepackage{hyperref}
\usepackage{fancyhdr}
\usepackage{titlesec}
\usepackage{tcolorbox}
\usepackage{qrcode}
\usepackage{graphicx}

% Build identity — same file as the passport and label. Fallbacks let this
% .tex compile standalone in a TeX editor without `make`.
\IfFileExists{build_info.tex}{\input{build_info}}{%
  \newcommand{\hwversion}{v2.1}%
  \newcommand{\fwversion}{dev}%
  \newcommand{\builddate}{unknown}%
  \newcommand{\buildcommit}{unknown}%
  \newcommand{\docrev}{unknown}%
}

% Public documentation site (GitHub Pages). Also the URL encoded in the
% QR code on the enclosure label — change in one place and everything follows.
\newcommand{\docsurl}{https://mrzlohex.github.io/dc-ups-2ch/}

\definecolor{accent}{HTML}{1E5B4F}
\definecolor{lightaccent}{HTML}{EAF3F1}
\definecolor{warn}{HTML}{8B1E1E}
\definecolor{lightwarn}{HTML}{F9ECEC}

\hypersetup{
  colorlinks=true,
  linkcolor=black,
  urlcolor=accent,
  pdfauthor={Изготовитель DC-UPS-2CH},
  pdftitle={Краткое руководство по эксплуатации DC-UPS-2CH \fwversion}
}

\setlength{\parindent}{0pt}
\setlength{\parskip}{0.4em}
\renewcommand{\arraystretch}{1.15}
\setlist[itemize]{leftmargin=5.5mm,itemsep=1.5pt,topsep=2pt}
\setlist[enumerate]{leftmargin=6.5mm,itemsep=2pt,topsep=2pt}

\titleformat{\section}{\large\bfseries\color{accent}}{\thesection.}{0.45em}{}
\titleformat{\subsection}{\normalsize\bfseries}{\thesubsection.}{0.45em}{}

\pagestyle{fancy}
\fancyhf{}
\lhead{\sffamily DC-UPS-2CH}
\rhead{\sffamily Руководство}
\cfoot{\sffamily\small \thepage}
\renewcommand{\headrulewidth}{0.35pt}
\setlength{\headheight}{13pt}

\newtcolorbox{infobox}[1][]{colback=lightaccent,colframe=accent,boxrule=0.7pt,arc=1.5mm,left=2mm,right=2mm,top=1.5mm,bottom=1.5mm,#1}
\newtcolorbox{warnbox}[1][]{colback=lightwarn,colframe=warn,boxrule=0.7pt,arc=1.5mm,left=2mm,right=2mm,top=1.5mm,bottom=1.5mm,#1}

\begin{document}

\begin{titlepage}
\centering
\vspace*{10mm}
{\sffamily\bfseries\fontsize{24}{28}\selectfont DC-UPS-2CH\\[2mm]}
{\sffamily\Large Краткое руководство по эксплуатации\\[5mm]}
{\sffamily\small для прошивки \fwversion\ (сборка \buildcommit, \builddate)}

\vspace{9mm}
\begin{infobox}
\centering
Двухканальный источник резервного питания для роутера, ONT и другого совместимого слаботочного оборудования.
\end{infobox}

\vspace{7mm}
\begin{tabularx}{\textwidth}{@{}lX@{}}
\toprule
Модель & DC-UPS-2CH \hwversion \\
Серийный номер & \rule{45mm}{0.4pt} \\
Дата передачи & \rule{45mm}{0.4pt} \\
Выход ROUTER & \rule{18mm}{0.4pt} В / \rule{18mm}{0.4pt} А \\
Выход ONT & \rule{18mm}{0.4pt} В / \rule{18mm}{0.4pt} А \\
\bottomrule
\end{tabularx}

\vfill
\begin{warnbox}
\textbf{Внимание.} Внутри корпуса присутствует сетевое напряжение 230~В. Пользователь не должен вскрывать корпус при подключённой сети или аккумуляторе.
\end{warnbox}
\vfill
{\small Хранить руководство рядом с устройством.}
\end{titlepage}

\section{Что делает устройство}

DC-UPS-2CH автоматически поддерживает питание двух независимых выходов. В обычном режиме устройство питается от сети и поддерживает аккумулятор в заряженном состоянии. При отключении 230~В переход на аккумулятор происходит автоматически, без действий пользователя.

\begin{tabularx}{\textwidth}{@{}p{44mm}X@{}}
\toprule
\textbf{Параметр} & \textbf{Значение} \\
\midrule
Сеть & 230~В AC, 50~Гц \\
Аккумулятор & 12~В, 7~А\,ч, AGM/SLA \\
Поддерживающий заряд & 13,8~В \\
Ограничение зарядного тока & 1,5~А \\
Предупреждение АКБ & 11,5~В \\
Отключение нагрузки (LVD) & 11,0~В \\
Разрешение повторного запуска & 12,5~В \\
\bottomrule
\end{tabularx}

\begin{infobox}[title=Нормальная работа]
При пропадании света ничего переключать не требуется. Выходы продолжают работу от аккумулятора, пока его напряжение не достигнет защитного порога.
\end{infobox}

\section{Ежедневное использование}

\subsection{Включение}
\begin{enumerate}
  \item Убедиться, что подключены нужные выходы.
  \item На корпусе есть отдельные механические выключатели выходов, перевести требуемые каналы во включённое положение.
  \item Подать 230~В.
  \item Дождаться запуска оборудования и Wi-Fi.
\end{enumerate}

В штатной эксплуатации дальнейшее ручное управление не требуется.

\subsection{При отключении электричества}
Устройство автоматически переходит на АКБ. При наличии связи отправляется уведомление ntfy. При снижении АКБ до 11,5~В отправляется предупреждение. При достижении 11,0~В оба управляемых выхода отключаются для защиты батареи.

После защитного отключения ESP32 переходит в аварийный глубокий сон. Она периодически просыпается и проверяет возврат внешнего питания. Когда сеть вернётся, устройство автоматически возобновит штатную работу.

\begin{warnbox}[title=Не путать с режимом хранения]
После аварийного глубокого сна устройство просыпается от возврата сети автоматически. В \textbf{режиме хранения} возврат 230~В устройство не будит.
\end{warnbox}

\section{Кнопка и индикация}

\subsection{Сервисная кнопка}
\begin{tabularx}{\textwidth}{@{}p{35mm}X@{}}
\toprule
\textbf{Действие} & \textbf{Результат} \\
\midrule
1 короткое нажатие & Отправить полный статус в ntfy \\
2 коротких нажатия & Включить \texttt{DC-UPS-Setup} примерно на 15 минут \\
Удержание 10~с & Перейти в режим хранения \\
Нажатие в режиме хранения& Разбудить устройство и выполнить обычный запуск \\
\bottomrule
\end{tabularx}

\subsection{Подтверждение кнопки}
\begin{itemize}
  \item 2 зелёных мигания --- статус успешно отправлен в ntfy;
  \item 3 красных мигания --- статус отправить не удалось;
  \item краткое совместное мигание LED --- запущен setup AP;
  \item 3 красных мигания перед выключением --- принят переход в режим хранения.
\end{itemize}

\subsection{Обычная индикация}
\begin{tabularx}{\textwidth}{@{}p{35mm}X@{}}
\toprule
Зелёный LED & наличие внешнего питания / сети \\
Красный LED & состояние АКБ, предупреждение или аварийный режим \\
\bottomrule
\end{tabularx}

\section{Режим хранения}

Режим хранения предназначен для перевозки или длительного хранения устройства, когда оба выхода и Wi-Fi должны быть полностью отключены.

\textbf{Чтобы выключить устройство для хранения:}
\begin{enumerate}
  \item Удерживать сервисную кнопку около 10 секунд.
  \item Дождаться трёх красных миганий.
  \item Отпустить кнопку.
  \item Выходы будут отключены, Wi-Fi выключится, ESP32 уйдёт в deep sleep без периодического таймера.
\end{enumerate}

\textbf{Чтобы снова включить устройство:} один раз нажать сервисную кнопку.

\begin{warnbox}[title=Важно]
В режими хранения устройство просыпается \textbf{только от физической кнопки}. Подключение или возврат 230~В само по себе его не включает.
\end{warnbox}

\section{Веб-панель и Wi-Fi}

При подключении к домашней сети панель доступна по адресу:
\begin{center}
\large\texttt{http://dc-ups.local/}
\end{center}

Также можно использовать IP-адрес, присланный через ntfy.

Через панель можно просматривать напряжения, состояние сети и выходов, Wi-Fi, интернет и журнал событий; управлять режимами ROUTER/ONT; выполнять restart каналов; менять Wi-Fi, ntfy и защитные параметры.

\subsection{Если домашний Wi-Fi недоступен}
Дважды быстро нажать сервисную кнопку. Появится точка доступа:

\begin{center}
\texttt{DC-UPS-Setup}
\end{center}

Пароль по умолчанию: \texttt{dc-ups-setup}, если он не был изменён при настройке.

\section{ntfy}

Устройство само отправляет уведомления о важных событиях. Кроме этого, в тот же topic можно отправить одну из команд:

\begin{center}
\texttt{!ups}\qquad \texttt{!ups status}\qquad \texttt{!ups ping}
\end{center}

Команды только запрашивают состояние и \textbf{не управляют питанием выходов}. Ответ включает напряжение АКБ, 24~В, состояние сети, ROUTER/ONT, интернет, uptime, RSSI и IP.

\begin{infobox}
Команды ntfy работают только когда ESP32 включена, подключена к Wi-Fi и имеет доступ в интернет. В аварийном глубоком сне и режиме хранения Wi-Fi выключен.
\end{infobox}

\section{Подключение оборудования}

Перед подключением нового роутера, ONT или другого устройства обязательно проверить:
\begin{itemize}
  \item требуемое выходное напряжение;
  \item полярность разъёма;
  \item допустимый ток соответствующего выхода;
  \item маркировку ROUTER / ONT на корпусе.
\end{itemize}

Для типового цилиндрического DC-разъёма данного экземпляра полярность должна соответствовать маркировке на корпусе. Не подключать оборудование, если напряжение или полярность неизвестны.

\section{Если что-то не работает}

\begin{tabularx}{\textwidth}{@{}p{43mm}X@{}}
\toprule
\textbf{Симптом} & \textbf{Что сделать} \\
\midrule
Нет питания на выходах & Проверить сеть, механические выключатели, состояние АКБ и режимы каналов в панели \\
Нет веб-панели & Попробовать \texttt{dc-ups.local}; затем дважды нажать кнопку и подключиться к setup AP \\
После отключения света выходы выключились & Возможен нормальный LVD при низком АКБ; дождаться возврата сети \\
После подачи 230~В устройство не запускается & Если ранее включён режим хранения  --- нажать сервисную кнопку \\
3 красных мигания после короткого нажатия & Статус в ntfy не отправлен; проверить Wi-Fi/интернет/topic \\
\bottomrule
\end{tabularx}

Если неисправность повторяется, не вскрывать устройство под напряжением и обратиться к изготовителю/сборщику.

\section{Безопасность и обслуживание}

\begin{itemize}
  \item не эксплуатировать устройство с открытым корпусом;
  \item не допускать воды, конденсата и перекрытия вентиляции;
  \item не замыкать клеммы аккумулятора;
  \item не превышать параметры выходов, указанные на наклейке;
  \item аккумулятор является расходным элементом и со временем теряет ёмкость;
  \item при заметном сокращении времени автономной работы аккумулятор следует проверить и при необходимости заменить в соответствии с полной сервисной документацией.
\end{itemize}

\begin{warnbox}
Перед вскрытием корпуса необходимо отключить 230~В и аккумулятор. Внутреннее обслуживание предназначено для подготовленного лица.
\end{warnbox}

\newpage
\section*{Быстрая памятка}

\begin{infobox}
\textbf{Обычная работа:} ничего переключать не нужно.\\
\textbf{Свет пропал:} устройство работает от АКБ автоматически.\\
\textbf{Свет вернулся:} после аварийного сна устройство восстановится автоматически.\\
\textbf{1x кнопка:} статус в ntfy.\\
\textbf{2x кнопка:} setup AP.\\
\textbf{10 с кнопка:} режим хранения.\\
\textbf{Из хранения:} нажать кнопку.\\
\textbf{Панель:} \texttt{http://dc-ups.local/}.\\
\textbf{ntfy:} \texttt{!ups status} или \texttt{!ups ping}.
\end{infobox}

\vspace{6mm}
\begin{center}
\begin{tabular}{@{}c@{\hspace{14mm}}c@{}}
\qrcode[height=28mm]{http://dc-ups.local/} & \qrcode[height=28mm]{\docsurl} \\[1.5mm]
\\
\\
{\small\parbox{40mm}{\centering Локальная веб-панель\\(та же LAN, поддержка mDNS)}} &
{\small\parbox{40mm}{\centering Онлайн-документация\\(полный паспорт, PDF)}} \\
\end{tabular}
\end{center}

\vfill
\begin{tabularx}{\textwidth}{@{}lX@{}}
Контакт изготовителя: & \rule{50mm}{0.4pt} \\
Дата документа: & \builddate \\
Редакция: & Руководство \fwversion\ / \docrev \\
\end{tabularx}

\end{document}
